% SEGMENT OLP-0415-S01
% Part: normal-modal-logic
% Chapter: syntax-and-semantics
% Section: modal-validity

\documentclass[../../../include/open-logic-section]{subfiles}

\begin{document}

\olfileid{nml}{syn}{val}

\olsection{செல்லுபடித்தன்மை}

\begin{explain}
  எல்லா மாதிரிகளிலும், அதாவது ஒவ்வொரு மாதிரியின் ஒவ்வோர் உலகிலும்,
  மெய்யாகும் !!^{formula}s குறிப்பாக ஆர்வமூட்டுகின்றன. சாத்தியமான
  உலகங்களின் மீதான ஏதோ ஓர் அணுகுறவால் பொருள்கோள் உருவாக்கப்படுகிறது என்ற
  பொருளில் அது ``இயல்பானதாக'' இருக்கும் வரையில், $\Box$, $\Diamond$
  ஆகியவை எவ்வாறு பொருள்கொள்ளப்பட்டாலும் மெய்யாகும் வாய்ப்புநிலைக்
  கூற்றுகளை அவை குறிக்கின்றன. இத்தகைய !!{formula}s ஐ
  \emph{செல்லுபடியானவை} என்கிறோம். எடுத்துக்காட்டாக, $\Box(p \land q)
  \lif \Box p$ செல்லுபடியாகும். எனினும், $\Box$ இன் மெய்மையியல்
  பொருள்கோளின் அடிப்படையில் செல்லுபடியாகும் என்று நாம் எதிர்பார்க்கக்கூடிய
  $\Box p \lif p$ போன்ற சில !!{formula}s செல்லுபடியானவை அல்ல. வெவ்வேறு
  வகையான அணுகுறவுகளைக் கொண்டு $\Box$, $\Diamond$ ஆகியவற்றின் வெவ்வேறு
  பொருள்கோள்களைப் பதிவு செய்ய முடிவது தொடர்புசார் மாதிரிகளின் ஈர்ப்பில்
  ஒரு பகுதியாகும். எனவே செல்லுபடித்தன்மையை \emph{எல்லா} மாதிரிகளுக்கும்
  சார்பாக மட்டுமன்றி, \emph{ஒரு குறிப்பிட்ட வகை மாதிரிகள் அனைத்துக்கும்}
  சார்பாக வரையறுக்க வேண்டும். எடுத்துக்காட்டாக, ஒவ்வோர் உலகையும் அதே
  உலகிலிருந்து அணுகக்கூடிய, அதாவது $R$ தற்சுட்டுப் பண்புடைய, எல்லா மாதிரிகளிலும்
  $\Box p \lif p$ மெய்யாகும் என்று தெரியவரும். மாதிரிகளின் வகுப்புகளுக்குச்
  சார்பாகச் செல்லுபடித்தன்மையை வரையறுப்பதால் இதனைச் சுருக்கமாகக் கூறலாம்:
  தற்சுட்டு மாதிரிகளின் வகுப்பில் $\Box p \lif p$ செல்லுபடியாகும்.
\end{explain}

\begin{defn}
  மாதிரிகளின் வகுப்பு $\mClass{C}$ இலுள்ள ஒவ்வொரு மாதிரியிலும்
  !!^a{formula}~$!A$ மெய்யாக இருந்தால் (அதாவது, $\mClass{C}$ இலுள்ள
  ஒவ்வொரு மாதிரியின் ஒவ்வோர் உலகிலும் மெய்யாக இருந்தால்), அது
  $\mClass{C}$ இல் \emph{செல்லுபடியாகும்}. $!A$ ஆனது $\mClass{C}$ இல்
  செல்லுபடியாக இருந்தால் $\mClass{C} \Entails !A$ என்று எழுதுகிறோம்;
  \emph{எல்லா} மாதிரிகளின் வகுப்பிலும் $!A$ செல்லுபடியாக இருந்தால்
  $\Entails !A$ என்று எழுதுகிறோம்.
\end{defn}

\begin{prop}\ollabel{prop:subset-class}
  $!A$ ஆனது $\mClass{C}$ இல் செல்லுபடியாக இருந்தால், ஒவ்வொரு வகுப்பு
  $\mClass{C}' \subseteq \mClass{C}$ இலும் அது செல்லுபடியாகும்.
\end{prop}

\begin{prop}\ollabel{prop:Nec-rule}
  $!A$ செல்லுபடியாக இருந்தால், $\Box!A$ உம் செல்லுபடியாகும்.
\end{prop}

\begin{proof}
  $\Entails !A$ என்க. $\Entails \Box!A$ என்று காட்ட, $\mModel{M} =
  \tuple{W, R, V}$ என்ற ஏதேனும் ஒரு மாதிரியையும் $w \in W$ என்ற ஏதேனும்
  ஓர் உலகையும் கொள்க. $Rww'$ எனில், $!A$ செல்லுபடியாக இருப்பதால்
  $\mSat{M}{!A}[w']$; எனவே $\mSat{M}{\Box!A}[w]$ உம். $\mModel{M}$ மற்றும்
  $w$ ஏதேனுமாக இருந்ததால், $\Entails \Box!A$.
\end{proof}

\begin{prob}
  பின்வருவன செல்லுபடியாகின்றன என்று காட்டுக:
  \begin{enumerate}
  \item $\Box p \lif \Box (q \lif p)$;
  \item $\Box \lnot \lfalse$;
  \item $\Box p \lif (\Box q \lif \Box p)$.
  \end{enumerate}
\end{prob}

\begin{prob}
  $W = \{w\}$ ஆகும் மாதிரிகள் $\mModel{M} = \tuple{W, R, V}$ இன் வகுப்பு
  $\mClass{C}$ இல் $!A \lif \Box!A$ செல்லுபடியாகும் என்று காட்டுக.
  அதேபோல், $R = \emptyset$ ஆகும் மாதிரிகள் $\mModel{M} = \tuple{W, R,
  V}$ இன் வகுப்பில் $!B \lif \Box !A$ மற்றும் $\Diamond !A \lif !B$
  செல்லுபடியாகின்றன என்று காட்டுக.
\end{prob}

\end{document}
